%! AP Statistics — Unit 9, Lesson 9.2 — Guided Notes
\documentclass[10pt]{article}
\usepackage{apstats-article}
\usepackage{apstats-boxes}

\begin{document}

\pageheader{Unit 9, Lesson 9.2}{Guided Notes}
\namedateperiod

\begin{objectivebox}
By the end of this lesson, I will be able to:
\begin{itemize}
    \item[$\square$] Identify a $t$-interval for slope as the appropriate procedure for estimating the true slope $\beta$.
    \item[$\square$] State the LINER conditions and verify each condition for a given context.
    \item[$\square$] Compute a confidence interval using $b \pm t^* \cdot SE_b$ with $df = n - 2$.
    \item[$\square$] Interpret a confidence interval for slope in context.
\end{itemize}
\end{objectivebox}

\begin{vocabbox}
\termblanklong{Population regression line ($\mu_y = \alpha + \beta x$)}
\termblanklong{Sample regression line ($\hat{y} = a + bx$)}
\termblanklong{Standard deviation of residuals ($s$)}
\termblanklong{Standard error of slope ($SE_b$)}
\termblanklong{$t$-interval for slope}
\termblanklong{LINER conditions}
\end{vocabbox}

\begin{hookbox}
A physical therapist randomly selects 25 patients and records the number of weekly PT
sessions ($x$) and recovery time in weeks ($y$). The LSRL is $\hat{y} = 12.4 - 0.8x$
with $s = 2.1$ and $SE_b = 0.22$. The therapist wants to estimate the true slope $\beta$
for all patients.

What procedure should she use, and why? \writelines{2}
\end{hookbox}

\begin{notesbox}{The Sampling Distribution of $b$ and the $t$-Interval (UNC-4.AC)}
The sample slope $b$ is a statistic that estimates the true slope $\beta$.
Its sampling distribution has:
\begin{itemize}
    \item Mean: $\mu_b = $ \blank{2cm}
    \item Standard deviation: $\sigma_b = \dfrac{\sigma}{\sigma_x\sqrt{n}}$
\end{itemize}

\vspace{0.1in}
Because $\sigma$ is unknown, we estimate it with
$s = $ \blank{7cm}.

\vspace{0.15in}
The appropriate confidence interval procedure for slope is a \blank{3cm} interval.

\vspace{0.1in}
\textbf{CI formula:}\quad
$b\ \pm\ t^* \cdot$ \blank{3cm},\quad where $df = n -$ \blank{1cm}.
\end{notesbox}

\begin{notesbox}{LINER Conditions for Inference on Slope (UNC-4.AD)}
\renewcommand{\arraystretch}{1.8}
\begin{tabular}{>{\bfseries}p{1.3cm} p{5.3cm} p{5.8cm}}
    Letter & Condition & How to Check \\
    \hline
    L &
        The true relationship between $x$ and $y$ is \blank{2.5cm}. &
        Residual plot shows \blank{4cm}. \\
    I &
        Observations are \blank{2.5cm}; if sampling w/o replacement: $n \le$ \blank{1.5cm}$\cdot N$. &
        Data come from a \blank{3.5cm}. \\
    N &
        The responses $y$ are approximately \blank{2.5cm} for fixed $x$. &
        Residual histogram; $n >$ \blank{1.5cm} if skewed. \\
    E &
        $\sigma_y$ does not \blank{2cm} with $x$ (equal variance). &
        Residual plot shows \blank{3.5cm}. \\
    R &
        Data come from a \blank{4cm} or randomized experiment. &
        Check study design. \\
\end{tabular}
\end{notesbox}

\begin{notesbox}{Computing a CI for Slope from Output (UNC-4.AE/AF)}
\textbf{Reading computer output:} Find $b$ in the \textbf{Coef} column and $SE_b$
in the \textbf{SE Coef} column on the slope row.

\vspace{0.1in}
Margin of error $= t^* \cdot$ \blank{3cm}.

\vspace{0.15in}
\textbf{Worked Example (PT Data):}\quad $b = -0.8$,\; $SE_b = 0.22$,\; $n = 25$,\; 95\% CI.

\vspace{0.05in}
$df = 25 - 2 = 23$;\quad $t^* = 2.069$

\vspace{0.1in}
$CI:\quad -0.8\ \pm\ 2.069 \times 0.22\ =\ -0.8\ \pm$ \blank{2cm}
$=\ ($ \blank{2.5cm},\; \blank{2.5cm}$)$

\vspace{0.15in}
\textbf{Interpret:} We are 95\% confident that the true slope of the population regression
line relating recovery time to weekly PT sessions for all patients is between
\blank{2cm} and \blank{2cm} weeks per session.
\end{notesbox}

\end{document}
